\documentclass[12pt]{article}
\usepackage[utf8]{inputenc}
\usepackage{amsmath, amssymb}
\usepackage{geometry}
\geometry{margin=1in}

\title{Problem Set: The Specific-Factors Model }
\author{International Trade Economics}
\date{}

\begin{document}
%\maketitle

\section*{Problem}

Canada produces two goods: Automobiles ($A$) and Wheat ($W$). Labor is mobile between the two sectors, while capital is specific to automobiles and land is specific to wheat.
\[
Q_A = K^{1/2}L_A^{1/2}, 
\qquad 
Q_W = T^{1/2}L_W^{1/2},
\qquad 
L_A + L_W = L.
\]

Endowments: $K = 49,\; T = 81,\; L = 120$ (so $K^{1/2}=7$ and $T^{1/2}=9$). \\
Prices: initially $P_A = 4,\; P_W = 2$. Later, $P_A$ rises to $6$ while $P_W$ remains $2$.

\bigskip

\noindent 1. Derive the marginal products of labor in each sector.

\bigskip

\noindent 2. Write the wage equalization condition and show the equation that determines the allocation of labor between the two sectors.

\bigskip

\noindent 3. Solve for the equilibrium allocation of labor $(L_A, L_W)$ and the common wage when $P_A=4$.

\bigskip

\noindent 4. Recalculate the equilibrium labor allocation and wage when $P_A=6$.

\bigskip

\noindent 6. Suppose the world price of automobiles is higher than Canada’s initial autarky price. Using the specific-factors model, explain the effects of opening to trade on: (a) labor allocation, (b) real returns to capital, land, and labor, and (c) Canada’s export/import pattern.

\bigskip
\hrule
\bigskip

\section*{Detailed Solutions}

\paragraph{Preliminaries.} With $Q_A=\sqrt{K}\sqrt{L_A}$ and $Q_W=\sqrt{T}\sqrt{L_W}$, the marginal products of labor are
\[
MPL_A=\frac{\partial Q_A}{\partial L_A}=\frac{1}{2}\,K^{1/2}L_A^{-1/2}=\frac{1}{2}\cdot 7\,L_A^{-1/2}=3.5\,L_A^{-1/2},
\]
\[
MPL_W=\frac{\partial Q_W}{\partial L_W}=\frac{1}{2}\,T^{1/2}L_W^{-1/2}=\frac{1}{2}\cdot 9\,L_W^{-1/2}=4.5\,L_W^{-1/2}.
\]
Labor market clearing: $L_W=120-L_A$.

\bigskip

\noindent \textbf{1. Marginal products.} As above,
\[
MPL_A=3.5\,L_A^{-1/2},\qquad MPL_W=4.5\,L_W^{-1/2}.
\]
Both decline in own labor (diminishing returns with the specific factor fixed).

\bigskip

\noindent \textbf{2. Wage equalization \& allocation equation.} With labor mobile,
\[
w=P_A\cdot MPL_A(L_A)=P_W\cdot MPL_W(L_W).
\]
Hence
\[
P_A\cdot 3.5\,L_A^{-1/2}=P_W\cdot 4.5\, (120-L_A)^{-1/2}.
\]
Given prices, this single equation in $L_A$ pins down the allocation; $L_W=120-L_A$ then follows, and the common wage is $w=P_A\cdot 3.5\,L_A^{-1/2}$ (or $w=P_W\cdot 4.5\,L_W^{-1/2}$).

\bigskip

\noindent \textbf{3. Initial equilibrium $(P_A=4,\;P_W=2)$.} Plug prices:
\[
4\cdot 3.5\,L_A^{-1/2}=2\cdot 4.5\,(120-L_A)^{-1/2}
\;\;\Longrightarrow\;\;
14\,L_A^{-1/2}=9\,(120-L_A)^{-1/2}.
\]
Square both sides (eliminates square roots):
\[
196\,L_A^{-1}=81\,(120-L_A)^{-1}.
\]
Cross-multiply:
\[
196(120-L_A)=81\,L_A.
\]
Solve for $L_A$:
\[
23520-196L_A=81L_A
\;\Rightarrow\;
23520=277L_A
\;\Rightarrow\;
\boxed{\,L_A=\frac{23520}{277}\approx 84.909\,}.
\]
Then $L_W=120-L_A\approx \boxed{35.091}$.

Common wage (either expression):
\[
w=14\,L_A^{-1/2}=\frac{14}{\sqrt{84.909}}\approx \boxed{1.519}.
\]

\emph{Outputs and specific-factor returns (useful for interpretation):}
\[
Q_A=\sqrt{K}\sqrt{L_A}=7\sqrt{84.909}\approx \boxed{64.53},\qquad
Q_W=\sqrt{T}\sqrt{L_W}=9\sqrt{35.091}\approx \boxed{53.32}.
\]
\[
r_K=P_A\cdot \frac{\partial Q_A}{\partial K}
=4\cdot \frac{1}{2}K^{-1/2}\sqrt{L_A}
=4\cdot \frac{\sqrt{L_A}}{14}
=\frac{2}{7}\sqrt{L_A}\approx \boxed{2.634},
\]
\[
r_T=P_W\cdot \frac{\partial Q_W}{\partial T}
=2\cdot \frac{1}{2}T^{-1/2}\sqrt{L_W}
=2\cdot \frac{\sqrt{L_W}}{18}
=\frac{\sqrt{L_W}}{9}\approx \boxed{0.658}.
\]

\bigskip

\noindent \textbf{4. New equilibrium $(P_A=6,\;P_W=2)$.} Update the wage condition:
\[
6\cdot 3.5\,L_A^{-1/2}=2\cdot 4.5\,(120-L_A)^{-1/2}
\;\;\Longrightarrow\;\;
21\,L_A^{-1/2}=9\,(120-L_A)^{-1/2}.
\]
Square:
\[
441\,L_A^{-1}=81\,(120-L_A)^{-1}.
\]
Cross-multiply:
\[
441(120-L_A)=81\,L_A
\;\Rightarrow\;
52920-441L_A=81L_A
\;\Rightarrow\;
52920=522L_A
\;\Rightarrow\;
\boxed{\,L_A=\frac{52920}{522}=\frac{2940}{29}\approx 101.379\,}.
\]
Then $L_W=120-L_A\approx \boxed{18.621}$.

Wage:
\[
w=21\,L_A^{-1/2}=\frac{21}{\sqrt{101.379}}\approx \boxed{2.086}.
\]

\emph{Specific-factor returns (for comparison):}
\[
r_K'=P_A'\cdot \frac{\partial Q_A}{\partial K}
=6\cdot \frac{\sqrt{L_A}}{14}
=\frac{3}{7}\sqrt{L_A}\approx \boxed{4.304},\qquad
r_T'=\frac{\sqrt{L_W}}{9}\approx \boxed{0.479}.
\]

\medskip
\emph{Comparative statics summary:} As $P_A$ rises from 4 to 5,
\[
L_A:\;84.9\rightarrow 101.4 \text{ (up)},\quad
w:\;1.519\rightarrow 2.086 \text{ (up but less than the }+50\%\text{ in }P_A),
\]
\[
r_K:\;2.63\rightarrow 4.30 \text{ (strongly up)},\quad
r_T:\;0.66\rightarrow 0.48 \text{ (down)}.
\]
Capital gains the most; land loses; labor’s nominal wage rises but not as much as $P_A$.

\bigskip

\noindent \textbf{5. Trade when world price of automobiles exceeds Canada’s autarky price.}

\emph{(a) Labor allocation.} With $P_A^{World}$ above the autarky $P_A$, the automobile wage schedule $w_A(L_A)=P_A^{World}\cdot 3.5\,L_A^{-1/2}$ shifts up relative to wheat. To restore $w_A(L_A)=w_W(L_W)$, labor moves \emph{into} automobiles: $L_A$ rises, $L_W=120-L_A$ falls.

\emph{(b) Factor returns.} 
- \textbf{Capital (specific to $A$):} Gains unambiguously. It benefits directly from the higher $P_A$ and indirectly from more labor in $A$ (which raises the marginal product of $K$). \\
- \textbf{Land (specific to $W$):} Loses. With less labor in wheat, the marginal product of land falls; $P_W$ is unchanged by assumption. \\
- \textbf{Labor (mobile):} Ambiguous in real terms. The common wage $w$ rises, but typically by \emph{less} than $P_A$, so $w/P_A$ falls while $w/P_W$ rises.

\emph{(c) Pattern of trade.} Canada exports \textbf{automobiles} and imports \textbf{wheat}. The higher world relative price of automobiles indicates Canada’s comparative advantage in $A$ at the margin; opening to trade expands $A$ output and contracts $W$.

\end{document}
